% Part: first-order-logic
% Chapter: syntax-and-semantics
% Section: first-order-languages

\documentclass[../../../include/open-logic-section]{subfiles}

\begin{document}

\olfileid{fol}{syn}{fol}

\olsection{د لومړۍ درجې ژبې}


د لومړۍ درجې د منطق عبارتونه له يوه بنسټيز لغت څخه جوړېږي چې
\emph{!!{variable}s}، \emph{!!{constant}s}، \emph{!!{predicate}s}
او کله ناکله \emph{!!{function}s} لري۔ له دغو سمبولونو سره، د منطقي
نښلوونکو، سورو او د وقف د نښو لکه قوسونو او کامو په کارولو،
\emph{ترمونه} او \emph{!!{formula}s} جوړېږي۔

\begin{explain}
په غير رسمي ډول، !!{predicate}s د ځانګړنو او اړيکو نومونه دي،
!!{constant}s د انفرادي څيزونو نومونه دي، او !!{function}s د نقشو
نومونه دي۔ دا ټول، له !!{identity}~$\eq$ پرته، \emph{غير منطقي
سمبولونه} دي او په ګډه يوه ژبه جوړوي۔ هره د لومړۍ درجې ژبه~$\Lang L$
د خپلو غير منطقي سمبولونو له مخې ټاکل کېږي۔ په تر ټولو عام حالت کښې،
$\Lang L$ د هر ډول سمبولونه په بې نهايته شمېر لري۔
\end{explain}

په عام حالت کښې موږ په لومړۍ درجې منطق کښې لاندې سمبولونه کاروو:

\begin{enumerate}
\item منطقي سمبولونه
\begin{enumerate}
\item منطقي نښلوونکي:
  \startycommalist
  \iftag{prvNot}{\ycomma $\lnot$ (نفي)}{}%
  \iftag{prvAnd}{\ycomma $\land$ (عطف)}{}%
  \iftag{prvOr}{\ycomma $\lor$ (فصل)}{}%
  \iftag{prvIf}{\ycomma $\lif$ (!!{conditional})}{}%
  \iftag{prvIff}{\ycomma $\liff$ (!!{biconditional})}{}%
  \iftag{prvAll}{\ycomma $\lforall$ (عمومي سور)}{}%
  \iftag{prvEx}{\ycomma $\lexists$ (وجودي سور)}{}.
\tagitem{prvFalse}{د !!{falsity}~$\lfalse$ لپاره بياني ثابت۔}{}
\tagitem{prvTrue}{د !!{truth}~$\ltrue$ لپاره بياني ثابت۔}{}
\item دوه ځايه !!{identity}~$\eq$۔
\item د !!{variable}s يو !!{denumerable}s سټ: $\Obj v_0$، $\Obj v_1$، $\Obj
  v_2$، \dots
\end{enumerate}
\item غير منطقي سمبولونه، چې د لومړۍ درجې منطق \emph{معياري
  ژبه} جوړوي
\begin{enumerate}
\item د هر $n>0$ لپاره د $n$-ځايه !!{predicate}s يو !!{denumerable}s سټ: $\Obj
  A^n_0$، $\Obj A^n_1$، $\Obj A^n_2$، \dots
\item د !!{constant}s يو !!{denumerable}s سټ: $\Obj c_0$، $\Obj c_1$، $\Obj
  c_2$، \dots۔
\item د هر $n>0$ لپاره د $n$-ځايه !!{function}s يو !!{denumerable}s سټ:
  $\Obj f^n_0$، $\Obj f^n_1$، $\Obj f^n_2$، \dots
\end{enumerate}
\item د وقف نښې: پرانيستے قوس (، تړلے قوس )، او کامه۔
\end{enumerate}

زموږ ډېری تعريفونه او پايلې به د لومړۍ درجې منطق د بشپړې معياري ژبې
لپاره بيانېږي۔ خو د کارونې له ډول سره سم ژبه يوازې څو !!{predicate}s،
!!{constant}s او !!{function}s ته هم محدودولے شو۔

\begin{ex}
د حساب ژبه~$\Lang L_A$ يو دوه ځايه !!{predicate}~$<$، يو
!!{constant}~$\Obj 0$، يو يو ځايه !!{function}~$\prime$، او دوه دوه
ځايه !!{function}s~$+$ او~$\times$ لري۔
\end{ex}

\begin{ex}
د سټ تيورۍ ژبه~$\Lang L_Z$ يوازې يو دوه ځايه !!{predicate}~$\in$
لري۔
\end{ex}

\begin{ex}
د ترتيبونو ژبه~$\Lang L_\le$ يوازې يو دوه ځايه !!{predicate}~$\le$
لري۔
\end{ex}

دا بيا هم يوازې دودونه دي: په رسمي ډول دا سمبولونه مستعار نومونه دي؛
د بېلګې په توګه، $<$، $\in$ او $\le$ د $\Obj A^2_0$، $\Obj 0$ د
$\Obj c_0$، $\prime$ د $\Obj f^1_0$، $+$ د $\Obj f^2_0$، او $\times$ د
$\Obj f^2_1$ مستعار نومونه دي۔

\iftag{defNot,defOr,defAnd,defIf,defIff,defTrue,defFalse,defEx,defAll}{%
پر هغو بنسټيزو نښلوونکو او
\iftag{notprvEx,notprvAll}{سور}{سورو} سربېره چې پورته معرفي شول، موږ
لاندې \emph{تعريف شوي} سمبولونه هم کاروو:
\startycommalist
  \iftag{defNot}{\ycomma $\lnot$ (نفي)}{}%
  \iftag{defAnd}{\ycomma $\land$ (عطف)}{}%
  \iftag{defOr}{\ycomma $\lor$ (فصل)}{}%
  \iftag{defIf}{\ycomma $\lif$ (!!{conditional})}{}%
  \iftag{defIff}{\ycomma $\liff$ (!!{biconditional})}{}%
  \iftag{defAll}{\ycomma $\lforall$ (عمومي سور)}{}%
  \iftag{defEx}{\ycomma $\lexists$ (وجودي سور)}{}%
  \iftag{defFalse}{\ycomma !!{falsity}~$\lfalse$}{}%
  \iftag{defTrue}{\ycomma !!{truth}~$\ltrue$}}{}.

\begin{tagblock}{defNot,defOr,defAnd,defIf,defIff,defTrue,defFalse,defEx,defAll}
\begin{explain}
تعريف شوے سمبول په رسمي ډول د ژبې برخه نۀ وي، بلکې د يوې غير رسمي
لنډونې په توګه ور پېژندل کېږي: د دې په مرسته هغه فارمولونه لنډولے شو
چې يوازې د بنسټيزو سمبولونو په کارولو ډېر اوږدېدل۔ دا يوه څرګنده ګټه
ده۔ تر دې لويه ګټه دا ده چې ثبوتونه لنډېږي۔ که يو سمبول بنسټيز وي، په
ثبوتونو کښې بايد جلا چلند ورسره وشي۔ نو څومره چې بنسټيز سمبولونه ډېر
وي، هومره زموږ ثبوتونه اوږدېږي۔
\end{explain}
\end{tagblock}

% Alternate symbols

\begin{intro}
کېدے شي تاسو له هغو اصطلاحاتو او سمبولونو سره بلد يئ چې زموږ له پورته
کارولو سره توپير لري۔ د منطق متنونه، او ښوونکي، عموماً د ``نفي''
دپاره $\sim$، $\neg$ او~!\، او د ``عطف'' دپاره $\wedge$، $\cdot$ او
$\&$ کاروي۔ د ``شرطي'' يا ``استلزام'' عام سمبولونه $\rightarrow$،
$\Rightarrow$ او $\supset$ دي۔
\iftag{prvIff,defIff}{د ``دوه اړخيز شرطي''، ``دوه اړخيز استلزام'' يا
  ``(مادي) هم ارزښت'' سمبولونه $\leftrightarrow$،
  $\Leftrightarrow$ او $\equiv$ دي۔}{}
\iftag{prvFalse,defFalse}{د $\lfalse$ سمبول ته په بېلو ځايونو کښې
  ``دروغوالے''، ``فالسوم''، ``محال'' يا ``ښکته'' ويل کېږي۔}{}
\iftag{prvTrue,defTrue}{د $\ltrue$ سمبول ته په بېلو ځايونو کښې
  ``رښتياوالے''، ``ويروم'' يا ``پورته'' ويل کېږي۔}{}

دود دا دے چې د لاتيني الفبا د پيل واړه توري (لکه $a$، $b$، $c$) د
!!{constant}s لپاره (چې کله ناکله نومونه بلل کېږي)، او د الفبا د پای
واړه توري (لکه $x$، $y$، $z$) د !!{variable}s لپاره وکارول شي۔ سورونه
له !!{variable}s سره يوځای راځي، لکه $x$؛ د ليکدود په بديلونو کښې د
عمومي سور لپاره $\forall x$، $(\forall x)$، $(x)$، $\Pi x$،
$\bigwedge_x$، او د وجودي سور لپاره $\exists x$، $(\exists
x)$، $(Ex)$، $\Sigma x$، $\bigvee_x$ شامل دي۔
\end{intro}

\begin{explain}
موږ ټول بياني عملګرونه او دواړه سورونه د ژبې بنسټيز سمبولونه ګڼلے شو۔
يا د بنسټيزو سمبولونو يوه کوچنۍ ټولګه غوره کولے او نور !!{operator}s
تعريف شوي ګڼلے شو۔ د بوليني عملګرونو په هغو ټولګو کښې چې ``د رښتياوالي
د تابع له مخې بشپړې'' دي، $\{ \lnot, \lor \}$، $\{ \lnot, \land \}$
او $\{ \lnot, \lif\}$ شامل دي---په بياني ځواک کښې د بشپړې لومړۍ
درجې ژبې دپاره له دې هره ټولګه له هر يوه سور سره يوځای کېدے شي۔

کېدے شي تاسو له دوو نورو !!{operator}s سره هم اشنا يئ: د شېفر
کرښه~$|$ (چې د هنري شېفر په نوم نومول شوې ده)، او د پيرس
غشی~$\downarrow$، چې د کواین خنجر هم بلل کېږي۔ کله چې دې عملګرونو ته په
ترتيب سره د ``نېنډ'' او ``نور'' معمولې لوستنې ورکړل شي، هر يو ئې په
يوازې ځان د رښتياوالي د تابع له مخې بشپړ وي۔
\end{explain}

\end{document}
